\documentclass[11pt]{amsart}
\usepackage{amsmath,amssymb,amsthm}
\usepackage[margin=1.05in]{geometry}

\numberwithin{equation}{section}
\newtheorem{theorem}{Theorem}[section]
\newtheorem{proposition}[theorem]{Proposition}
\newtheorem{lemma}[theorem]{Lemma}
\newtheorem{corollary}[theorem]{Corollary}
\theoremstyle{definition}
\newtheorem{definition}[theorem]{Definition}
\theoremstyle{remark}
\newtheorem{remark}[theorem]{Remark}
\newtheorem{example}[theorem]{Example}

\newcommand{\Z}{\mathbb{Z}}
\newcommand{\Q}{\mathbb{Q}}
\newcommand{\R}{\mathbb{R}}
\newcommand{\N}{\mathbb{N}}
\newcommand{\C}{\mathbb{C}}
\newcommand{\cA}{\mathcal{A}}
\newcommand{\cG}{\mathcal{G}}
\newcommand{\cR}{\mathcal{R}}
\newcommand{\cF}{\mathcal{F}}
\newcommand{\pp}{\mathfrak{p}}
\newcommand{\aaa}{\mathfrak{a}}
\newcommand{\Ns}{\mathrm{N}}
\newcommand{\Prob}{\operatorname{Prob}}
\newcommand{\Out}{\operatorname{Out}}
\newcommand{\Aut}{\operatorname{Aut}}
\newcommand{\Msym}{M_{\mathrm{sym}}}
\newcommand{\phisym}{\varphi_{\mathrm{sym}}}
\newcommand{\id}{\mathrm{id}}

\PassOptionsToPackage{hyphens}{url}
\usepackage[hidelinks,bookmarks=false]{hyperref}

\title[Amenability and the failure of superrigidity]{Localized Bost--Connes factors are amenable:\\ the failure of $W^*$-superrigidity,\\ and what survives at the $C^*$-level}
\author{Ruqing Chen}
\address{GUT Geoservice Inc., Montreal, Canada}
\email{ruqing@hotmail.com}
\thanks{Sources, code and data for the entire series: \url{https://github.com/Ruqing1963/localized-bost-connes}; this paper is in \texttt{papers/IX-rigidity/}.}
\date{}

\begin{document}

\begin{abstract}
We settle the rigidity question for the type $\mathrm{III}$ factors of the localized
Bost--Connes systems, and the answer is negative in the strongest available sense.

The groupoid $\cG_S=I_S\ltimes\widetilde\Omega_S$ has abelian, hence amenable, acting group
$I_S=\bigoplus_{\pp\in S}\Z$, so $M_{K,S,\beta}$ is an injective factor. By the classification of
injective factors (Connes, Krieger, Haagerup) its isomorphism class is determined by its flow
of weights alone --- by nothing beyond the single number $\lambda$ in type
$\mathrm{III}_\lambda$, by nothing at all in type $\mathrm{III}_1$. The group
$\Sigma_\beta=\mathrm{pr}_\R(\Xi^{\mathrm{ext}}_\beta)$ of \cite{SeriesI} bounds that flow, lying
in its point spectrum, and determines it whenever the spectrum is pure point; whatever
arithmetic the factor retains is therefore in the flow of weights, and in the cases decided
in \cite{SeriesI} there is none.

Consequently \emph{class field $W^*$-superrigidity is false, maximally}. For $S=P_K$ and
$0<\beta\le1$ every number field gives type $\mathrm{III}_1$ \cite[Thm.~3.2]{SeriesI}, hence
\[
M_{\Q,P_\Q,1}\ \cong\ M_{\Q,P_\Q,1/2}\ \cong\ M_{K,P_K,\beta}\ \cong\ R_\infty
\]
for every number field $K$ and every $\beta\in(0,1]$: the factor remembers neither $K$, nor
$S$, nor $\beta$. Popa's deformation/rigidity machinery cannot be applied, and not for
technical reasons: $s$-malleable deformations produce rigidity through a spectral gap, which
requires the acting group to be non-amenable, and $I_S$ is abelian. Cartan uniqueness holds,
but as a consequence of Connes--Feldman--Weiss rather than of Popa, and only up to conjugacy
by an automorphism --- not up to unitary conjugacy --- while orbit equivalence rigidity fails
outright, all amenable relations of a given type being orbit equivalent by Krieger.

What survives is at the $C^*$-level and in inclusions. The obstruction group $\Xi_\beta$ is
an invariant of the $C^*$-dynamical system $(\cA_{K,S},\sigma)$ through its
$\mathrm{KMS}_\beta$ simplex, and is destroyed by taking the weak closure in a single
extremal state; this is the operator-algebraic form of the measure-class principle of
Paper III of the series. And the symmetry-breaking inclusion
$\Msym^{G}\subset\Msym$ of Paper II has index $|\Xi_\beta|$, so the \emph{pair} recovers the
order of the obstruction group even though neither algebra alone recovers anything.
\end{abstract}

\maketitle

\tableofcontents

\section{Amenability, and what it costs}

Throughout, $K$ is a number field, $S$ an infinite set of finite primes, $\beta>0$ with
$\zeta_{K,S}(\beta)=\infty$, $\varphi$ an extremal $\mathrm{KMS}_\beta$ state and
\[
M_{K,S,\beta}=\pi_\varphi\bigl(\cA_{K,S}\bigr)''=W^*\bigl(\cG_S,\nu\bigr),
\qquad \cG_S=I_S\ltimes\widetilde\Omega_S,\quad I_S=\bigoplus_{\pp\in S}\Z .
\]
\emph{The series.} The foundational paper is \cite{Main} and Paper I is \cite{SeriesI};
Papers II, III, V, VII and XXVII, cited by numeral, are in preparation, and nothing proved
here depends on them.


\begin{lemma}[Amenability]\label{lem:amenable}
$\cG_S$ is a second countable étale groupoid with abelian, hence amenable, acting group;
therefore $\cG_S$ is amenable, $C^*_r(\cG_S)=C^*(\cG_S)$, and $M_{K,S,\beta}$ is an injective
von Neumann algebra. It is a factor of type $\mathrm{III}$ by \cite[Prop.~6.1]{Main}.
\end{lemma}

\begin{proof}
$I_S$ is a countable abelian group, hence amenable, and a transformation groupoid by an
amenable group acting on a locally compact space is amenable; amenability of the groupoid
gives injectivity of the associated von Neumann algebra, equivalently the orbit relation
of $I_S$ on $(\widetilde\Omega_S,\nu)$ is amenable and its von Neumann algebra
\cite{FM} is injective \cite{CFW}.
\end{proof}

\begin{theorem}[The complete $W^*$-invariant]\label{thm:complete}
$M_{K,S,\beta}$ is an injective factor of type $\mathrm{III}$, so its isomorphism class is
determined by its flow of weights, and by nothing else:
\begin{itemize}
\item if the type is $\mathrm{III}_1$, then $M\cong R_\infty$, the unique injective factor of
type $\mathrm{III}_1$ \cite{Haagerup};
\item if the type is $\mathrm{III}_\lambda$, $0<\lambda<1$, then $M\cong R_\lambda$, the Powers
factor \cite{Connes76};
\item if the type is $\mathrm{III}_0$, then $M$ is the Krieger factor of its flow of weights,
which is a complete invariant \cite{Krieger}.
\end{itemize}
The group $\Sigma_\beta=\mathrm{pr}_\R(\Xi^{\mathrm{ext}}_\beta)$ of \cite[Def.~1.1]{SeriesI}
constrains the flow of weights: $\Sigma_\beta$ is contained in the point spectrum $T(M)$ of
the flow \cite[Rem.~2.3]{SeriesI}, so $\Sigma_\beta\ne\{0\}$ excludes $\mathrm{III}_1$,
$\Sigma_\beta$ dense forces $\mathrm{III}_0$, and $\Sigma_\beta=c\Z$ forces
$\mathrm{III}_{e^{-2\pi k/c}}$ for some $k\ge1$ or $\mathrm{III}_0$ \cite[Prop.~2.2]{SeriesI}.
If the flow of weights has pure point spectrum, $\Sigma_\beta$ determines it and hence the
factor; whether it always does is open. In the two cases decided in \cite{SeriesI}, namely
$S=P_K$ with $0<\beta\le1$ (type $\mathrm{III}_1$, \cite[Thm.~3.2]{SeriesI}) and the band set
of \cite[Thm.~4.2]{SeriesI} (type $\mathrm{III}_{e^{-1}}$), the factor is $R_\infty$ and
$R_{e^{-1}}$ respectively.
\end{theorem}

\begin{proof}
Lemma~\ref{lem:amenable} and the cited classification theorems: Connes classified the
injective factors of type $\mathrm{III}_\lambda$, $\lambda\ne1$, the $\mathrm{III}_0$ case
being reduced to Krieger's theorem on the flow of weights, and Haagerup settled
$\mathrm{III}_1$. The constraints are \cite[Prop.~2.2 and Rem.~2.3]{SeriesI}; the two
decided cases are \cite[Thm.~3.2, Thm.~4.2]{SeriesI}.
\end{proof}

\begin{remark}[What the factor can still remember]\label{rem:remember}
The one place where arithmetic could survive the weak closure is the flow of weights in
type $\mathrm{III}_0$, or the number $\lambda$ in type $\mathrm{III}_\lambda$. Both are
functions of the essential-value data of the Radon--Nikodym cocycle $\Ns\aaa^{-\beta}$, hence
of the sequence of norms together with the measure class of $\pi_\beta$; by
\cite[Prop.~6.1]{Main} the Connes invariant satisfies $S(M)\subseteq\overline{\Gamma_\beta}\cup\{0\}$
with $\Gamma_\beta$ the group generated by the $\Ns\pp^{-\beta}$. What the factor remembers
is therefore at most an ergodic flow built from the norm data and the measure class; two
systems with the same flow of weights are isomorphic whatever their norms, fields and
temperatures, so none of these is reconstructible from the factor.
\end{remark}

\section{$W^*$-superrigidity fails maximally}

\begin{theorem}[No reconstruction]\label{thm:fails}
Let $K_1,K_2$ be any number fields and $\beta_1,\beta_2\in(0,1]$. Then
\[
M_{K_1,P_{K_1},\beta_1}\ \cong\ M_{K_2,P_{K_2},\beta_2}\ \cong\ R_\infty .
\]
In particular the isomorphism class of the factor determines neither the number field $K$,
nor the prime set $S$, nor the inverse temperature $\beta$, nor the norms $\{\Ns\pp\}$, nor
the symmetry group $G_S$, nor the obstruction group $\Xi_\beta$.
\end{theorem}

\begin{proof}
By \cite[Thm.~3.2]{SeriesI} the type is $\mathrm{III}_1$ for $S=P_K$ and $0<\beta\le1$, for
every number field $K$; by Theorem~\ref{thm:complete} the factor is then the unique injective
$\mathrm{III}_1$ factor.
\end{proof}

\begin{example}[Explicit collapses]\label{ex:collapse}
\[
\begin{array}{lllll}
K & S & \beta & \text{type} & M\\\hline
\Q & P_\Q & 1 & \mathrm{III}_1 & R_\infty\\
\Q & P_\Q & 1/2 & \mathrm{III}_1 & R_\infty\\
\Q(i) & P_K & 1 & \mathrm{III}_1 & R_\infty\\
\Q(2^{1/3}) & P_K & 0.9 & \mathrm{III}_1 & R_\infty\\
\Q & \text{band set of \cite[Thm.~4.2]{SeriesI}} & 1 & \mathrm{III}_{e^{-1}} & R_{e^{-1}}
\end{array}
\]
The first four are isomorphic von Neumann algebras although their arithmetic input differs in
every respect. The fifth is distinguishable from them, but only through the single number
$e^{-1}$, and it has the \emph{same} factor as every other system of type
$\mathrm{III}_{e^{-1}}$ --- for instance the system of any $K$ and $S$ whose norms are all
powers of a single $\lambda_0$ with $\lambda_0^{-\beta}=e^{-1}$, when it is of type
$\mathrm{III}_{e^{-1}}$ rather than $\mathrm{III}_0$ \cite[Prop.~6.1]{Main}.
\end{example}

\begin{corollary}[The point spectrum is not an invariant]\label{cor:pointspec}
The set $\{\Ns\pp^{-\beta}:\pp\in S\}$ is not a $W^*$-invariant. Only the closed subgroup
$\overline{\Gamma_\beta}$ enters Connes' invariant $S(M)$, by \cite[Prop.~6.1]{Main}, and
$S(M)$ determines only the type. In particular no reconstruction of the prime norms, and
hence no norm-preserving bijection $S_1\cong S_2$, can be extracted from an isomorphism of
factors.
\end{corollary}

\begin{proof}
The four systems of Example~\ref{ex:collapse} with $S=P_K$ have pairwise different norm
sets and the same factor. In general $S(M)=\bigcap_{\psi}\mathrm{Sp}(\Delta_\psi)$ is a closed
subset of $[0,\infty)$ determined by the type, and $S(M)\subseteq\overline{\Gamma_\beta}\cup\{0\}$
\cite[Prop.~6.1]{Main}; norm sets generating the same closed subgroup $\overline{\Gamma_\beta}$
are indistinguishable by $S(M)$, and by Theorem~\ref{thm:complete} the factor carries only
the flow of weights.
\end{proof}

\section{Why Popa's machinery does not apply}

\begin{proposition}[No spectral gap]\label{prop:nogap}
Let $(\alpha_t)_{t\in\R}$ be the $s$-malleable Bernoulli deformation associated with the
product measure $\pi_\beta=\bigotimes_\pp\mu_{t_\pp}$ on $V=\Z_{\ge0}^S$. The deformation
exists, but the conclusion of Popa's spectral gap argument is unavailable: that argument
derives rigidity from the failure of the acting group to have almost invariant vectors, and
$I_S$ is abelian, hence amenable, hence has almost invariant vectors in every
representation. No subgroup of $I_S$ has property $(\mathrm T)$ or a spectral gap.
\end{proposition}

\begin{proof}
An abelian group is amenable and the regular representation of an infinite amenable group has
almost invariant vectors; Popa's deformation/rigidity arguments \cite{Popa} require the
opposite, a rigid subalgebra or a spectral gap for the deformation.
\end{proof}

\begin{remark}[This is not a gap in the exposition]\label{rem:notgap}
Deformation/rigidity theory exists to produce rigidity where amenability fails. Applied to an
amenable algebra it can only recover what Connes' classification already gives, which by
Theorem~\ref{thm:complete} is the type. The absence of a Popa-type theorem here is therefore
structural, and the correct reading of Theorem~\ref{thm:fails} is not that the rigidity proof
is missing but that the rigidity statement is false.
\end{remark}

\section{Cartan subalgebras and orbit equivalence}

\begin{theorem}[Cartan uniqueness, correctly attributed]\label{thm:cartan}
$A=L^\infty(Y_S,\nu)\subset M_{K,S,\beta}$ is a Cartan subalgebra, and any two Cartan
subalgebras of $M_{K,S,\beta}$ are conjugate by an \emph{automorphism} of $M$. This is a
consequence of Connes--Feldman--Weiss, not of Popa intertwining, and it does \emph{not} give
unitary conjugacy: since $\Out(M)$ is nontrivial for the injective factors, conjugacy by an
automorphism is strictly weaker.
\end{theorem}

\begin{proof}
$A$ is maximal abelian with normalizer generating $M$, by the groupoid picture, the groupoid
being principal on a conull set \cite[Prop.~2.8]{Main}. A Cartan subalgebra $B\subset M$ gives
by Feldman--Moore \cite{FM} an ergodic non-singular relation $\cR_B$ and a $2$-cocycle with
$(M,B)\cong(W^*(\cR_B,\sigma_B),L^\infty)$; $\cR_B$ is amenable because $M$ is injective, hence
hyperfinite by Connes--Feldman--Weiss \cite{CFW}, and for hyperfinite relations the
$2$-cocycle is a coboundary \cite{FM}. Two Cartan subalgebras $B_1,B_2$ thus give two ergodic
hyperfinite relations whose associated factors coincide, hence with the same flow of weights;
by Krieger's theorem \cite{Krieger} they are orbit equivalent, so the Cartan pairs $(M,B_1)$
and $(M,B_2)$ are isomorphic, i.e.\ conjugate by an automorphism of $M$. Unique Cartan
decomposition up to unitary conjugacy is a rigidity phenomenon of non-amenable factors
\cite{Houdayer}, and does not hold here.
\end{proof}

\begin{corollary}[Orbit equivalence rigidity fails]\label{cor:oe}
An isomorphism $M_{K_1,S_1,\beta_1}\cong M_{K_2,S_2,\beta_2}$ does induce an isomorphism of
the underlying orbit equivalence relations, by Theorem~\ref{thm:cartan} and Feldman--Moore.
But this carries no information: by Krieger's theorem all amenable ergodic non-singular
relations of a given type are already orbit equivalent, so the induced isomorphism exists
whether or not the arithmetic inputs are related.
\end{corollary}

\begin{remark}[The Radon--Nikodym cocycle]\label{rem:rn}
The cocycle $D(v+a,g;v,g)=\Ns\aaa^{-\beta}$ is preserved modulo coboundaries by any orbit
equivalence, and its essential range is the ratio set. But the ratio set is
$S(M)\cap(0,\infty)$, which by Theorem~\ref{thm:complete} is the type. The cocycle therefore
transmits exactly the information already counted, and no more; in particular it does not
transmit the point spectrum, by Corollary~\ref{cor:pointspec}.
\end{remark}

\section{What survives}

\begin{theorem}[$\Xi_\beta$ is a $C^*$-invariant, not a $W^*$-invariant]\label{thm:Cstar}
The obstruction group $\Xi_\beta(K,S)$ is determined by the $C^*$-dynamical system
$(\cA_{K,S},\sigma)$, through its $\mathrm{KMS}_\beta$ simplex
$\Prob(G_S/\Xi_\beta^\perp)$ \cite[Thm.~3.6]{Main}. It is \emph{not} determined by
$M_{K,S,\beta}$, by Theorem~\ref{thm:fails}. The passage from the $C^*$-algebra with its
flow to the weak closure in a single extremal state destroys it.
\end{theorem}

\begin{proof}
The first statement is the classification theorem; the second is Theorem~\ref{thm:fails},
since two systems with $\Xi_{\beta_1}\ne\Xi_{\beta_2}$ can have isomorphic factors.
\end{proof}

\begin{remark}[Consistency with the measure-class principle]\label{rem:measureclass}
Theorem~\ref{thm:Cstar} is the operator-algebraic form of the principle of Paper III of the
series, where $\Xi_\beta$ is identified as an invariant of the measure class of $\pi_\beta$,
all the $\beta$-dependence sitting there. Taking the weak closure in a single extremal state retains
the measure class only through its Krieger invariant, that is through the type. The present
theorem is therefore not a new phenomenon but the same one, seen at the level of von Neumann
algebras rather than of cohomology.
\end{remark}

\begin{theorem}[The inclusion remembers the order]\label{thm:inclusion}
Let $\Xi_\beta$ be finite, let $\Msym=\pi_{\phisym}(\cA_{K,S})''=\bigoplus_{\varepsilon}M_\varepsilon$
be the algebra of the symmetric state, the sum over the $|\Xi_\beta|$ extremal
$\mathrm{KMS}_\beta$ states, and let $\Msym^{G}\subset\Msym$ be the fixed-point algebra of the
action of $G=G_S/\Xi_\beta^\perp$, which permutes the summands simply transitively and
identifies $\Msym^G$ with any one $M_\varepsilon$ embedded diagonally. Then
\[
\bigl[\Msym:\Msym^{G}\bigr]_0=\bigl|\Xi_\beta\bigr| ,
\]
the minimal index, so the isomorphism class of the \emph{pair} determines $|\Xi_\beta|$,
even though by Theorem~\ref{thm:fails} neither $M_\varepsilon$ nor $\Msym$ need determine
anything. Equivalently, $\log[\Msym:\Msym^G]_0=\log|\Xi_\beta|$ equals the relative entropy
of an extremal phase against the symmetric one.
\end{theorem}

\begin{proof}
The action of $G_S$ on the $\mathrm{KMS}_\beta$ simplex factors through $G$ and permutes the
extreme points simply transitively \cite[Thm.~3.6(3)]{Main}, hence the summands of $\Msym$,
whose GNS representations are pairwise disjoint; so $\Msym^G$ is the diagonal copy
$D\otimes1\subseteq D\otimes\ell^\infty(G)$ with $D=M_\varepsilon$, whose minimal index is
$|G|$, realized by the averaging expectation. The relative entropy of one extremal state
against the uniform mixture of $|G|$ mutually singular states is $\log|G|$. These are the
computations of Paper II of the series (its Theorem 5.8 and Corollary 5.11), reproduced
here for completeness.
\end{proof}

\begin{remark}[Where to look for rigidity]\label{rem:whererigidity}
Theorem~\ref{thm:inclusion} indicates the level at which rigidity questions are not
vacuous: not the single factor, which is $R_\infty$, $R_\lambda$ or a Krieger factor, but
structured objects --- the $C^*$-algebra with its flow, the inclusion $\Msym^G\subset\Msym$,
the Cartan pair together with the $G_S$-action. Each of these carries strictly more than the factor, and the
first carries $\Xi_\beta$ in full.
\end{remark}

\section{Automorphisms, and the fundamental group}

\begin{proposition}[Both invariants are degenerate]\label{prop:degenerate}
\begin{enumerate}
\item The fundamental group $\cF(M)$ is not defined for type $\mathrm{III}$ factors, being an
invariant of $\mathrm{II}_1$ factors; the analogous statement, that $M\otimes B(H)\cong M$,
holds trivially for every properly infinite factor.
\item $\Out(M_{K,S,\beta})$ is that of the injective factor determined by
Theorem~\ref{thm:complete} --- of $R_\infty$ or $R_\lambda$ --- and so is the same for all
systems of the same type, carrying no arithmetic information.
\end{enumerate}
\end{proposition}

\begin{proof}
(1) is the definition. (2) is Theorem~\ref{thm:complete}: isomorphic algebras have isomorphic
outer automorphism groups.
\end{proof}

\begin{remark}[On the proposed embedding $G_S/\Xi_\beta^\perp\hookrightarrow\Out(M)$]\label{rem:out}
The group $G_S$ acts on the $\mathrm{KMS}_\beta$ simplex, permuting the extremal states
simply transitively through $G_S/\Xi_\beta^\perp$ \cite[Thm.~3.6(3)]{Main}. It therefore acts
on $\Msym$, and gives isomorphisms between the different summands $M_\varepsilon$; but those
summands are mutually isomorphic by Theorem~\ref{thm:complete}, so the resulting
automorphisms of $M_\varepsilon$ are not canonically defined and depend on the chosen
identifications. We do not claim an embedding $G_S/\Xi_\beta^\perp\hookrightarrow\Out(M_\varepsilon)$
and we do not see how to obtain one; what is canonical is the action on the pair, which is
the content of Theorem~\ref{thm:inclusion}.
\end{remark}

\section{Assessment}

\[
\begin{array}{lll}
\text{Statement} & \text{Status} & \text{Reference}\\\hline
M_{K,S,\beta}\ \text{is injective} & \text{true} & \text{Lem.~1.1}\\
\text{complete }W^*\text{-invariant}=\text{flow of weights} & \text{true; }\Sigma_\beta\text{ bounds it} & \text{Thm.~1.2}\\
\text{class field }W^*\text{-superrigidity} & \textbf{false, maximally} & \text{Thm.~2.1}\\
\text{reconstruction of }\beta & \text{false} & \text{Thm.~2.1}\\
\text{reconstruction of }\{\Ns\pp\} & \text{false} & \text{Cor.~2.3}\\
\text{reconstruction of }G_S,\ \Xi_\beta & \text{false} & \text{Thm.~2.1}\\
s\text{-malleable deformation gives rigidity} & \text{false: no spectral gap} & \text{Prop.~3.1}\\
\text{Cartan unique up to unitary conjugacy} & \text{false} & \text{Thm.~4.1}\\
\text{Cartan unique up to automorphism} & \text{true, by CFW} & \text{Thm.~4.1}\\
\text{OE rigidity} & \text{false, by Krieger} & \text{Cor.~4.2}\\
\Xi_\beta\ \text{a }C^*\text{-invariant} & \textbf{true} & \text{Thm.~5.1}\\
{[\Msym:\Msym^G]_0}=|\Xi_\beta| & \textbf{true} & \text{Thm.~5.3}\\
\cF(M),\ \Out(M)\ \text{carry arithmetic} & \text{false} & \text{Prop.~6.1}
\end{array}
\]

The picture is uniform with the rest of this project. In Paper III of the series the
arithmetic was shown to be an invariant of a measure class, invisible to $K$-theory, to groupoid homology
and to dynamical entropy, all of which discard the Radon--Nikodym derivative. Passing to the
von Neumann algebra discards it too, retaining only its Krieger invariant; and since the
groupoid is amenable, Connes' classification then leaves nothing at all. The failure recorded
here is not a failure of technique. It is the same measure-class principle, applied one level
further out.

Three questions. First, is there a non-amenable enlargement of the system to which
deformation/rigidity does apply? Paper V of the series shows that a free non-abelian
semigroup empties the high-temperature region, and Paper XXVII that non-abelian symmetry
can be obtained with the abelian group $I_S$ acting through a non-abelian cocycle --- but
the groupoid is then still amenable, the factor still injective, and nothing here changes;
an Ore semigroup that is non-amenable but retains nested range projections would be the
object to look for, and we do not know whether one exists. Second, how much does the Cartan pair together with the
$G_S$-action remember? Theorem~\ref{thm:inclusion} recovers $|\Xi_\beta|$ from the inclusion;
whether the triple $(M,A,G_S)$ recovers $\Xi_\beta$ itself, or even $S$, is open. Third, at
the $C^*$-level the classification of \cite{Main} is complete, but we have not asked whether
$\cA_{K,S}$ itself, without its flow, remembers anything; by Paper II its $K$-theory does
not, and the question of $C^*$-isomorphism is untouched.

\begin{thebibliography}{9}
\bibitem{Connes76} A. Connes, \emph{Classification of injective factors. Cases $\mathrm{II}_1$, $\mathrm{II}_\infty$, $\mathrm{III}_\lambda$, $\lambda\ne1$}, Ann. of Math. (2) 104 (1976), 73--115.
\bibitem{CFW} A. Connes, J. Feldman, B. Weiss, \emph{An amenable equivalence relation is generated by a single transformation}, Ergodic Theory Dynam. Systems 1 (1981), 431--450.
\bibitem{FM} J. Feldman, C. C. Moore, \emph{Ergodic equivalence relations, cohomology, and von Neumann algebras II}, Trans. Amer. Math. Soc. 234 (1977), 325--359.
\bibitem{Haagerup} U. Haagerup, \emph{Connes' bicentralizer problem and uniqueness of the injective factor of type $\mathrm{III}_1$}, Acta Math. 158 (1987), 95--148.
\bibitem{Houdayer} C. Houdayer, S. Vaes, \emph{Type $\mathrm{III}$ factors with unique Cartan decomposition}, J. Math. Pures Appl. (9) 100 (2013), 564--590.
\bibitem{Krieger} W. Krieger, \emph{On ergodic flows and the isomorphism of factors}, Math. Ann. 223 (1976), 19--70.
\bibitem{Popa} S. Popa, \emph{Strong rigidity of $\mathrm{II}_1$ factors arising from malleable actions of $w$-rigid groups I}, Invent. Math. 165 (2006), 369--408.
\bibitem{Main} R. Chen, \emph{KMS state uniqueness and phase transitions in localized Bost--Connes systems: from Chebotarev density to the Hardy--Littlewood conjecture}, preprint, 2026, DOI \href{https://doi.org/10.5281/zenodo.22152101}{10.5281/zenodo.22152101}.
\bibitem{SeriesI} R. Chen, \emph{Localized Bost--Connes systems I: factor type and the spectrum of transitions}, preprint, 2026, DOI \href{https://doi.org/10.5281/zenodo.22160827}{10.5281/zenodo.22160827}.
\end{thebibliography}

\end{document}
